\documentclass[romanian,12pt,a4paper]{article}

\title{Generator procedural de obiecte 3D}
\author{Andrei N. Onea}
\date{Februarie 2023}

\usepackage[utf8]{inputenc}
\usepackage[T1]{fontenc}
\usepackage{babel}
\usepackage{array}
\usepackage{graphicx}
\usepackage{multirow}
\usepackage{caption}
\usepackage{subcaption}
\usepackage{xcolor}
\usepackage{listings}
\usepackage[toc]{appendix}
\usepackage{multirow}
\usepackage[sorting=none]{biblatex}
\usepackage{csquotes}
\addbibresource{src/bib/main.bib}

\linespread{1.2}

\begin{document}

	\newcolumntype{P}[1]{>{\centering\arraybackslash}p{#1}}

	\begin{titlepage}

		\centering
		\vspace*{-2cm}
		\includegraphics[width=0.5\textwidth]{src/img/logo.png}\par
		\vspace{0.15cm}
		{\scshape Facultatea de Automatică și Calculatoare \par}
		\vspace{3cm}
		{\huge\scshape Generator procedural de obiecte 3D \par}
		\vspace{1cm}
		{\scshape Lucrare de masterat \par}
		\vspace{0.15cm}
		{rev 0.1 \par}
		\vfill

		\makebox[\textwidth][c]{
			\centering
			\begin{tabular}{P{0.5\textwidth} P{0.5\textwidth}}
				\textit{Coordonator}
				&
				\textit{Autor}
				\\
				Dr. Vice Decan Alin \scshape Moldoveanu
				&
				Andrei-Nicolae \scshape Onea
			\end{tabular}
		}
		\vfill

		{\small Februarie 2023 \par}

	\end{titlepage}

	\begin{abstract}
		Ceva informație aici...
	\end{abstract}

	\pagebreak

	\section*{Mulțumiri}

		Aș vrea să mulțumesc în primul rând părinților mei pentru tot sprijinul și sfaturile lor, dar mai ales pentru că m-au suportat în perioada aceasta de studii. Nu au învățat nimic din precedentele...

	\pagebreak

	\tableofcontents

	\clearpage

	\printbibliography

\end{document}